\chapter{Química computacional: Hartree--Fock y DFT}\label{ch:b3:computational-chemistry}

Se cree que el ion hidruro de helio, \ce{HeH+}, fue la primera molécula que se
formó al enfriarse el Universo joven: un átomo de helio que retiene un protón. Se
obtuvo en un tubo de descarga de laboratorio en 1925, pero se buscó en el espacio durante
décadas sin éxito, hasta que en 2019 se detectó su raya de rotación en una
nebulosa planetaria. Para entonces, su longitud de enlace, sus frecuencias de
vibración y las longitudes de onda de sus rayas se conocían desde hacía mucho ---por
cálculo. Un programa informático que resuelve, de forma aproximada, la \omterm{thm:b3:quantum-model-systems:schrodinger}{ecuación de
Schrödinger} de los electrones puede predecir la forma, la energía
y el espectro de una molécula que nadie ha metido en un frasco. Este capítulo
explica cómo funcionan esos programas, desde la separación de núcleos y electrones
hasta la teoría de Hartree--Fock y la \omterm{def:b3:computational-chemistry:dft}{teoría del funcional de la densidad}, y para qué
se puede confiar en sus números y para qué no.

\begin{recall}
El \cref{ch:b3:quantum-model-systems} aportó el hamiltoniano, la \omterm{thm:b3:quantum-model-systems:schrodinger}{ecuación de
Schrödinger} y el átomo de hidrógeno; el \cref{ch:b3:many-electron-atoms}, el \omterm{def:b3:many-electron-atoms:slater}{determinante
de Slater} y la \omterm{def:b3:many-electron-atoms:exchange}{integral de intercambio}. El volumen del segundo año construyó orbitales
moleculares como combinaciones lineales de orbitales atómicos (CLOA), con integrales de
solapamiento, de Coulomb y de resonancia y un determinante secular, y resolvió el
método de Hückel; también usó los espectros fotoelectrónicos para ver las energías orbitales.
\end{recall}

\begin{omfigure}
\includegraphics[width=0.6\linewidth]{images/book4/ai/b3-03-cluster.jpg}

{\small Un clúster de cálculo: filas de servidores en los que se ejecutan hoy la mayoría de los
cálculos de química cuántica, a menudo durante horas o días por molécula.}
\end{omfigure}

\section{Born--Oppenheimer y la superficie de energía potencial}

El hamiltoniano de una molécula contiene las energías cinéticas de sus núcleos y
electrones y todas sus interacciones coulombianas. Los núcleos son al menos 1836
veces más pesados que los electrones: para las mismas fuerzas se mueven mucho más
despacio, y los electrones se ajustan casi instantáneamente a cada posición de los
núcleos.

\begin{definition}[Aproximación de Born--Oppenheimer]\label{def:b3:computational-chemistry:born-oppenheimer}
En la \emph{aproximación de Born--Oppenheimer}\index{aproximacion de Born--Oppenheimer@aproximación de Born--Oppenheimer},
la \omterm{thm:b3:quantum-model-systems:schrodinger}{ecuación de Schrödinger} electrónica se resuelve con los núcleos
fijos en las posiciones $\mathbf R$,
\[
\hat H_{\mathrm{el}}\,\psi_{\mathrm{el}}(\mathbf r;\mathbf R) =
E_{\mathrm{el}}(\mathbf R)\,\psi_{\mathrm{el}}(\mathbf r;\mathbf R),
\]
y los núcleos se mueven después en el potencial $E_{\mathrm{el}}(\mathbf R) +
V_{\mathrm{nn}}(\mathbf R)$, la energía electrónica más la repulsión de los
núcleos.
\end{definition}

\begin{definition}[Superficie de energía potencial, punto estacionario, optimización de geometría]\label{def:b3:computational-chemistry:pes}
La \emph{superficie de energía potencial}\index{superficie de energia potencial@superficie de energía potencial} (SEP) de
una molécula es la función $U(\mathbf R) = E_{\mathrm{el}}(\mathbf R) +
V_{\mathrm{nn}}(\mathbf R)$ de sus coordenadas nucleares. Un \emph{punto
estacionario}\index{punto estacionario} es un punto en el que se anulan todas las derivadas
primeras de $U$; un mínimo es una \emph{geometría de equilibrio}\index{geometria de equilibrio@geometría de equilibrio}.
Buscar uno a partir de una estructura de partida, siguiendo las
fuerzas $-\nabla U$, es una \emph{optimización de geometría}\index{optimizacion de geometria@optimización de geometría}.
\end{definition}

Para una molécula diatómica, la SEP es una curva $U(R)$; para una triatómica, una
función de tres coordenadas; para $N$ átomos, de $3N - 6$. La masa de los
núcleos no interviene en $\hat H_{\mathrm{el}}$: \ce{H2}, \ce{HD} y \ce{D2} comparten la
misma SEP, la misma longitud de enlace y la misma \omterm{def:b3:quantum-model-systems:oscillator}{constante de fuerza}, y solo difieren en cómo
se mueven los núcleos sobre ella (\omterm{def:b3:quantum-model-systems:oscillator}{energía de punto cero}, \cref{ch:b3:quantum-model-systems}).

\begin{proposition}[Frecuencias a partir de la hessiana]\label{prop:b3:computational-chemistry:hessian}
Cerca de un mínimo, $U \approx U_0 + \frac12\sum_{ij}H_{ij}\,\delta q_i\,\delta q_j$, con
la hessiana $H_{ij} = \partial^2U/\partial q_i\partial q_j$. En coordenadas
ponderadas por la masa, $\delta q_i\sqrt{m_i}$, los \omterm{def:b3:quantum-model-systems:operator}{valores propios} $\lambda_k$ de la
hessiana ponderada por la masa dan las frecuencias angulares armónicas $\omega_k =
\sqrt{\lambda_k}$ de los modos normales. En un mínimo, todos los $\lambda_k$ son
positivos (seis, o cinco para una molécula lineal, son nulos: traslaciones y
rotaciones); en un punto de silla de primer orden, exactamente uno es negativo.
\end{proposition}

\begin{proof}[Demostración parcial]
Para una diatómica, $U \approx U_0 + \frac12k(R - R_e)^2$ con $k = U''(R_e)$, y
el oscilador del \cref{ch:b3:quantum-model-systems} da $\omega =
\sqrt{k/\mu}$. En general, las ecuaciones clásicas $m_i\ddot q_i =
-\sum_jH_{ij}q_j$ se convierten, en coordenadas ponderadas por la masa $x_i = \sqrt{m_i}q_i$,
en $\ddot{\mathbf x} = -\tilde H\mathbf x$ con la matriz simétrica $\tilde H_{ij} =
H_{ij}/\sqrt{m_im_j}$; sus vectores propios oscilan independientemente a
$\sqrt{\lambda_k}$. Un $\lambda_k$ negativo da una frecuencia imaginaria: en
esa dirección la energía baja; es un máximo. Los modos normales del
\cref{ch:b3:group-theory-applied} son estos vectores propios.
\end{proof}

\begin{omfigure}
\begin{tikzpicture}[font=\small, scale=0.95]
  \fill[omDef!8] (0,0) ellipse (1.9 and 1.15);
  \foreach \s in {1.0,0.75,0.5,0.25}{\draw[omDef!70] (0,0) ellipse ({1.9*\s} and {1.15*\s});}
  \fill[omThm!8] (6,0.4) ellipse (1.6 and 1.0);
  \foreach \s in {1.0,0.75,0.5,0.25}{\draw[omThm!70] (6,0.4) ellipse ({1.6*\s} and {1.0*\s});}
  \draw[gray] (2.1,-0.15) .. controls (2.9,0.55) and (3.4,0.55) .. (4.3,0.35);
  \draw[gray] (2.2,0.9) .. controls (3.0,0.35) and (3.5,0.35) .. (4.1,1.2);
  \node[omsaddle] at (3.15,0.42) {};
  \draw[ompath] (0,0) .. controls (1.8,0.6) and (2.4,0.5) .. (3.15,0.42) .. controls (4.0,0.35) and (4.8,0.4) .. (6,0.4);
  \node[anchor=north] at (0,-1.25) {mínimo A (reactivos)};
  \node[anchor=north] at (6,-0.7) {mínimo B (productos)};
  \node[anchor=south] at (3.15,1.25) {punto de silla};
  \draw[->, gray] (3.15,1.25) -- (3.15,0.6);
  \fill (0,0) circle (1.2pt) (6,0.4) circle (1.2pt);
  \draw[->] (-2.3,-2.2) -- (8.2,-2.2) node[right] {$q_1$};
  \draw[->] (-2.3,-2.2) -- (-2.3,1.8) node[above] {$q_2$};
\end{tikzpicture}

{\small Una \omterm{def:b3:computational-chemistry:pes}{superficie de energía potencial} vista desde arriba, como un mapa de curvas de nivel sobre dos
coordenadas nucleares (esquema). Dos mínimos están unidos, a través de un punto de
silla (cruz), por el camino de menor energía (trazo discontinuo): el camino de energía mínima
estudiado en el \cref{ch:b3:rate-theories}.}
\end{omfigure}

\section{El principio variacional}

La \omterm{thm:b3:quantum-model-systems:schrodinger}{ecuación de Schrödinger} electrónica no puede resolverse exactamente para ninguna
molécula con más de un electrón. Las aproximaciones se juzgan con un solo
teorema.

\begin{theorem}[El principio variacional]\label{thm:b3:computational-chemistry:variational}
Para toda función de prueba normalizable $\phi$ que cumpla las condiciones de
contorno del problema,
\[
E[\phi] = \frac{\langle\phi|\hat H|\phi\rangle}{\langle\phi|\phi\rangle} \ge E_0,
\]
donde $E_0$ es el \omterm{def:b3:quantum-model-systems:operator}{valor propio} más bajo de $\hat H$; la igualdad solo se da si $\phi$
es una \omterm{def:b3:quantum-model-systems:operator}{función propia} del estado fundamental. Es el \emph{principio
variacional}\index{principio variacional}.
\end{theorem}

\begin{proof}
Se desarrolla $\phi = \sum_nc_n\psi_n$ sobre las \omterm{def:b3:quantum-model-systems:operator}{funciones propias} ortonormales de $\hat H$
(\cref{thm:b3:quantum-model-systems:hermitian-real}). Entonces $\langle\phi|\hat
H|\phi\rangle = \sum_n|c_n|^2E_n \ge E_0\sum_n|c_n|^2 = E_0\langle\phi|\phi\rangle$,
ya que todo $E_n \ge E_0$; la igualdad exige $c_n = 0$ siempre que $E_n > E_0$.
\end{proof}

Cuanto menor es la energía de una función de prueba, mejor es; todo parámetro
ajustable se fija minimizando $E$.

\begin{proposition}[Helio con una carga apantallada]\label{prop:b3:computational-chemistry:helium-zeta}
Para el helio, la función de prueba $\phi = \eu^{-\zeta(r_1 + r_2)}$ (en unidades atómicas:
longitudes en $a_0$, energías en $E_h$) da $E(\zeta) = \zeta^2 - \frac{27}{8}\zeta$,
mínima para $\zeta = \frac{27}{16}$, con $E = -(\frac{27}{16})^2 E_h =
\qty{-2.8477}{\hartree}$.
\end{proposition}

\begin{proof}
Para un electrón en $\eu^{-\zeta r}$, $\langle T\rangle = \zeta^2/2$ y
$\langle 1/r\rangle = \zeta$; la repulsión de dos electrones en la misma función $1s$
es $\frac58\zeta$ (una integral clásica, que se admite). Así, $E = 2\cdot\frac12
\zeta^2 - 2\cdot2\zeta + \frac58\zeta = \zeta^2 - \frac{27}{8}\zeta$, y $\dd E/\dd
\zeta = 0$ da $\zeta = 27/16 = 1.6875$.
\end{proof}

\begin{example}[¿Cuán buena es?]\label{ex:b3:computational-chemistry:helium}
% ledger: ie:He, ie4:He2, const:Eh, const:eV
La energía total del helio es menos la suma de sus dos energías de ionización,
$-(24.5874 + 54.4178)\,\unit{eV} = \qty{-2.9034}{\hartree}$. La función de un
parámetro queda \qty{0.056}{\hartree} (\qty{1.5}{eV}) demasiado alta: cada electrón ve un
núcleo apantallado hasta $Z_{\mathrm{eff}} = 1.69$ por el otro, como sugieren
las reglas de Slater, pero los electrones también se esquivan de un modo que ningún producto de dos
funciones puede describir.
\end{example}

\begin{proposition}[Variación lineal]\label{prop:b3:computational-chemistry:linear-variation}
Para una función de prueba $\phi = \sum_{i=1}^n c_i\chi_i$ construida sobre $n$ funciones fijas,
con $H_{ij} = \langle\chi_i|\hat H|\chi_j\rangle$ y $S_{ij} = \langle\chi_i|\chi_j
\rangle$, los valores estacionarios de $E$ son las raíces de $\det(H - ES) = 0$, y
la raíz más baja es una cota superior de $E_0$.
\end{proposition}

\begin{proof}
$E\sum_{ij}c_ic_jS_{ij} = \sum_{ij}c_ic_jH_{ij}$ (coeficientes reales). Se deriva
respecto de $c_k$ y se impone $\partial E/\partial c_k = 0$: $\sum_j(H_{kj} -
ES_{kj})c_j = 0$ para todo $k$. Una solución no nula exige que se anule el
determinante. La raíz más baja es $E[\phi]$ para su vector propio, luego $\ge E_0$.
\end{proof}

El método de Hückel del volumen del segundo año es esta proposición con orbitales
$p$ y $H_{ij}$ empíricos. Todos los métodos siguientes son la misma idea con
mejores funciones y mejores elementos de matriz.

\section{Conjuntos de base}

\begin{definition}[Conjunto de base, base mínima, orbitales de tipo Slater y de tipo gaussiano]\label{def:b3:computational-chemistry:basis}
Un \emph{conjunto de base}\index{conjunto de base} es el conjunto de funciones fijas $\chi_i$,
centradas en los átomos, a partir de las cuales se construyen los orbitales moleculares. Una
\emph{base mínima}\index{base minima@base mínima} tiene una función por cada
orbital atómico ocupado (una $1s$ para H, cinco para C). Un \emph{orbital de tipo
Slater}\index{orbital de tipo Slater} (STO) tiene la forma radial
$r^{n-1}\eu^{-\zeta r}$; un \emph{orbital de tipo gaussiano}\index{orbital de tipo gaussiano}
(GTO) tiene la forma $\eu^{-\alpha r^2}$ multiplicada por un polinomio en $x$, $y$,
$z$. Una \emph{función gaussiana contraída}\index{funcion gaussiana contraida@función gaussiana contraída}
es una combinación fija $\sum_kd_k\,\eu^{-\alpha_kr^2}$ de gaussianas primitivas.
\end{definition}

Las funciones de Slater tienen la forma correcta (un punto anguloso en el núcleo y una
cola exponencial), pero sus integrales bielectrónicas sobre varios centros son
muy costosas; las funciones gaussianas tienen la forma equivocada, pero el producto de dos
gaussianas centradas en puntos distintos es una gaussiana centrada en un punto entre ambos, de modo que
toda integral se reduce a una fórmula cerrada. Las contracciones obtienen lo mejor de
ambas: en la base STO-3G, cada función de Slater se sustituye por una suma fija de
tres gaussianas ajustadas a ella.

\begin{example}[La función STO-3G del hidrógeno]\label{ex:b3:computational-chemistry:sto3g}
% ledger: bse:H.sto3g.a1, bse:H.sto3g.a2, bse:H.sto3g.a3, bse:sto3g.c1, bse:sto3g.c2, bse:sto3g.c3
La función $1s$ del hidrógeno en STO-3G es $0.15433\,g(3.42525) + 0.53533\,
g(0.62391) + 0.44463\,g(0.16886)$, donde $g(\alpha)$ es la gaussiana normalizada
$(2\alpha/\pi)^{3/4}\eu^{-\alpha r^2}$ (exponentes en $a_0^{-2}$). Imita una
función de Slater de exponente $\zeta = 1.24$ ---un átomo de hidrógeno ligeramente
comprimido, como lo está en las moléculas. Su solapamiento con la función $1s$ exacta
del mismo exponente es $0.9998$; solo se pierden el punto anguloso en el núcleo y la
cola lejana (figura siguiente).
\end{example}

\begin{definition}[Funciones de valencia desdoblada, de polarización y difusas]\label{def:b3:computational-chemistry:split-valence}
Una \emph{base de valencia desdoblada}\index{base de valencia desdoblada} describe cada
orbital de valencia mediante dos (o más) funciones de tamaños distintos, de modo que el
orbital pueda crecer o encogerse en una molécula. Una \emph{función de
polarización}\index{funcion de polarizacion@función de polarización} tiene un $l$ mayor que los orbitales
ocupados ($p$ en H, $d$ en C) y permite que la densidad se desplace fuera del átomo; una
\emph{función difusa}\index{funcion difusa@función difusa} tiene un exponente pequeño y
describe aniones e interacciones débiles lejos de los núcleos.
\end{definition}

Un nombre como 6-31G(d) dice: seis primitivas para cada función de núcleo,
la valencia desdoblada en una contracción de tres y una sola primitiva, y \omterm{def:b3:computational-chemistry:split-valence}{funciones de polarización}
$d$ en los átomos pesados. Los resultados convergen hacia el
\emph{límite de base completa} a medida que crece la base; un cálculo solo vale lo que
vale su base.

\begin{omfigure}
\begin{tikzpicture}
\begin{axis}[width=0.46\linewidth, height=5.4cm, xmin=0, xmax=5, ymin=0,
  ymax=0.6, axis lines=left, xlabel={$r$ / $a_0$}, ylabel={$\phi(r)$},
  legend cell align=left, legend style={font=\small, draw=none},
  tick label style={font=\small}, label style={font=\small},
  title={\small STO-3G frente a $\eu^{-r}$ ($\zeta = 1$)}]
\addplot[black, thick] table[x=r, y=slater] {figdata/out/bachelor-3-minimal-basis-hf-sto.dat};
\addlegendentry{Slater}
\addplot[omDef, thick, dashed] table[x=r, y=g3] {figdata/out/bachelor-3-minimal-basis-hf-sto.dat};
\addlegendentry{STO-3G}
\addplot[omThm, thick, dotted] table[x=r, y=g1] {figdata/out/bachelor-3-minimal-basis-hf-sto.dat};
\addlegendentry{una gaussiana}
\end{axis}
\end{tikzpicture}\hfill
\begin{tikzpicture}
\begin{axis}[width=0.46\linewidth, height=5.4cm, xmin=0.5, xmax=8.5,
  ymin=-2.86, ymax=-2.78, axis lines=left, xlabel={ciclo SCF},
  ylabel={$E$ / $E_h$}, ytick={-2.85,-2.83,-2.81,-2.79},
  yticklabel style={/pgf/number format/fixed, /pgf/number format/precision=2},
  tick label style={font=\small}, label style={font=\small},
  title={\small SCF del \ce{HeH+}}]
\addplot[omDef, thick, mark=*] table[x=it, y=E] {figdata/out/bachelor-3-minimal-basis-hf-heh-scf.dat};
\addplot[gray, dashed, domain=0.5:8.5] {-2.84184};
\end{axis}
\end{tikzpicture}

{\small Izquierda: la contracción STO-3G (trazo discontinuo) reproduce de cerca la función de Slater
$1s$ salvo en el núcleo; una sola gaussiana (punteado) no lo hace.
Derecha: la energía de \ce{HeH+} a \qty{1.4632}{\bohr} durante las iteraciones
de \omterm{def:b3:computational-chemistry:hartree-fock}{campo autoconsistente} (STO-3G); converge por arriba hacia
\qty{-2.8418}{\hartree} en pocos ciclos, y cada iteración obedece al \omterm{thm:b3:computational-chemistry:variational}{principio
variacional}.}
\end{omfigure}

\section{Teoría de Hartree--Fock}

La \omterm{def:b3:many-electron-atoms:indistinguishable}{función de onda antisimétrica} más sencilla de $N$ electrones es un único \omterm{def:b3:many-electron-atoms:slater}{determinante
de Slater}. La teoría de Hartree--Fock encuentra el mejor.

\begin{definition}[Método de Hartree--Fock, operador de Fock, campo autoconsistente]\label{def:b3:computational-chemistry:hartree-fock}
El \emph{método de Hartree--Fock}\index{metodo de Hartree--Fock@método de Hartree--Fock} aproxima el
estado fundamental por el \omterm{def:b3:many-electron-atoms:slater}{determinante de Slater} único de menor energía. Sus orbitales
son \omterm{def:b3:quantum-model-systems:operator}{funciones propias} del \emph{operador de Fock}\index{operador de Fock}
\[
\hat f(1) = \hat h(1) + \sum_b\big[2\hat J_b(1) - \hat K_b(1)\big],
\]
en el caso de capa cerrada: $\hat h$ es la energía cinética y la atracción hacia los
núcleos de un electrón, $\hat J_b$ la repulsión de la nube de carga del
orbital $b$, y $\hat K_b$ el \omterm{def:b3:quantum-model-systems:operator}{operador} de intercambio, cuyos \omterm{def:b3:quantum-model-systems:hermitian}{valores esperados} son
las \omterm{def:b3:many-electron-atoms:exchange}{integrales de intercambio} del \cref{ch:b3:many-electron-atoms}. Como $\hat f$
depende de los orbitales que determina, las ecuaciones se resuelven por
iteración hasta que los orbitales dejan de cambiar: un \emph{campo
autoconsistente}\index{campo autoconsistente} (SCF).
\end{definition}

\begin{theorem}[Ecuaciones de Roothaan--Hall]\label{thm:b3:computational-chemistry:roothaan}
Escribiendo cada orbital en una base, $\phi_a = \sum_\mu C_{\mu a}\chi_\mu$, las
ecuaciones de Hartree--Fock de capa cerrada se convierten en el problema matricial de \omterm{def:b3:quantum-model-systems:operator}{valores propios}
\[
\mathbf F\mathbf C = \mathbf S\mathbf C\boldsymbol\varepsilon, \qquad F_{\mu\nu} =
h_{\mu\nu} + \sum_{\lambda\sigma}P_{\lambda\sigma}\big[(\mu\nu|\sigma\lambda) -
\tfrac12(\mu\lambda|\sigma\nu)\big],
\]
con la matriz de densidad $P_{\lambda\sigma} = 2\sum_a^{\mathrm{occ}}C_{\lambda a}C_{\sigma
a}$ y las integrales bielectrónicas $(\mu\nu|\lambda\sigma) = \iint\chi_\mu(1)\chi_\nu(1)
r_{12}^{-1}\chi_\lambda(2)\chi_\sigma(2)$ (unidades atómicas).
\end{theorem}

\begin{proof}[Demostración parcial]
La energía del determinante es $E = \sum_{\mu\nu}P_{\mu\nu}h_{\mu\nu} + \frac12
\sum P_{\mu\nu}P_{\lambda\sigma}[(\mu\nu|\sigma\lambda) - \frac12(\mu\lambda|\sigma\nu)]$, una
función cuadrática de los coeficientes. Se minimiza con las restricciones
de que los orbitales sigan siendo ortonormales, $\mathbf C^{\mathsf T}\mathbf S\mathbf C = \mathbf 1$,
con un multiplicador de Lagrange para cada restricción: las condiciones de
estacionariedad son $\mathbf F\mathbf C = \mathbf S\mathbf C\boldsymbol\varepsilon$, y los
multiplicadores forman una matriz que puede diagonalizarse mediante una rotación de los
orbitales ocupados entre sí. El álgebra se admite.
\end{proof}

\begin{method}[El procedimiento SCF]\label{met:b3:computational-chemistry:scf}
\begin{enumerate}
  \item Calcular una sola vez las integrales $S_{\mu\nu}$, $h_{\mu\nu}$ y $(\mu\nu|\lambda\sigma)$.
  \item Proponer una matriz de densidad de partida (por ejemplo, a partir de los vectores propios de
        $\mathbf h$ sola, la estimación ``de núcleo'').
  \item Construir $\mathbf F$ a partir de $\mathbf P$; resolver $\mathbf F\mathbf C = \mathbf S\mathbf C
        \boldsymbol\varepsilon$ (ortogonalizar con $\mathbf S^{-1/2}$ y después diagonalizar).
  \item Llenar los orbitales más bajos, con dos electrones cada uno; formar la nueva $\mathbf P$
        y la energía.
  \item Repetir desde el paso 3 hasta que la energía y $\mathbf P$ cambien menos que un
        umbral.
\end{enumerate}
\end{method}

\begin{omfigure}
\begin{tikzpicture}[font=\small, node distance=4mm,
  box/.style={draw=black!60, rounded corners=2pt, fill=omDef!6, align=center,
  text width=4.4cm, inner sep=4pt}]
  \node[box] (a) {integrales $S$, $h$, $(\mu\nu|\lambda\sigma)$};
  \node[box, below=of a] (b) {$\mathbf P$ de partida};
  \node[box, below=of b] (c) {construir $\mathbf F(\mathbf P)$};
  \node[box, below=of c] (d) {resolver $\mathbf F\mathbf C = \mathbf S\mathbf C\boldsymbol\varepsilon$};
  \node[box, below=of d] (e) {nueva $\mathbf P$, energía $E$};
  \node[diamond, aspect=2.4, draw=black!60, fill=omThm!6, below=of e, inner sep=1pt] (f) {¿converge?};
  \node[box, below=of f, fill=omProp!8] (g) {orbitales, $\varepsilon$, $E$};
  \draw[->] (a) -- (b); \draw[->] (b) -- (c); \draw[->] (c) -- (d);
  \draw[->] (d) -- (e); \draw[->] (e) -- (f);
  \draw[->] (f) -- node[left] {sí} (g);
  \draw[->] (f.east) -- ++(1.6,0) |- node[pos=0.25, right] {no} (c.east);
\end{tikzpicture}

{\small El bucle de \omterm{def:b3:computational-chemistry:hartree-fock}{campo autoconsistente} de un cálculo de Hartree--Fock.}
\end{omfigure}

\begin{example}[\ce{H2} en la base mínima]\label{ex:b3:computational-chemistry:h2}
% ledger: re:H2, janaf:H.dfH0, diat:H2.we, diat:H2.wexe
Con una función STO-3G por átomo, el orbital enlazante queda fijado por la simetría,
$\sigma_g \propto \chi_A + \chi_B$, y el SCF necesita un solo paso. A $R =
\qty{1.4}{\bohr}$, el solapamiento es $S_{AB} = 0.6593$ y la energía
\qty{-1.1167}{\hartree}; el mínimo está a \qty{1.346}{\bohr} (\qty{71.2}{pm},
frente a los \qty{74.1}{pm} medidos) y \qty{-1.1175}{\hartree}. La figura siguiente
compara la curva con la verdadera, construida a partir de datos espectroscópicos
medidos y de la energía de disociación (una curva de Morse, \cref{ch:b3:rovibrational-spectroscopy}).
\end{example}

\begin{omfigure}
\begin{tikzpicture}
\begin{axis}[width=0.7\linewidth, height=6cm, xmin=0.8, xmax=10, ymin=-1.2,
  ymax=-0.45, axis lines=left, xlabel={$R$ / $a_0$}, ylabel={$E$ / $E_h$},
  legend cell align=left, legend style={font=\small, at={(0.04,0.98)}, anchor=north west, draw=none},
  tick label style={font=\small}, label style={font=\small}]
\addplot[omDef, thick] table[x=R, y=rhf] {figdata/out/bachelor-3-minimal-basis-hf-h2.dat};
\addlegendentry{RHF/STO-3G}
\addplot[black, thick, dashed] table[x=R, y=morse] {figdata/out/bachelor-3-minimal-basis-hf-h2.dat};
\addlegendentry{medida (forma de Morse)}
\addplot[gray, dotted, domain=0.8:10] {-1.0};
\addplot[omThm, dotted, domain=0.8:10] {-0.93316};
\node[font=\scriptsize, anchor=north] at (axis cs:8,-1.005) {dos átomos H exactos};
\node[font=\scriptsize, anchor=south, omThm] at (axis cs:8,-0.930) {dos átomos H en STO-3G};
\end{axis}
\end{tikzpicture}

{\small La energía de \ce{H2} en función de la distancia internuclear: Hartree--Fock
restringido en la base STO-3G (trazo continuo) y la curva verdadera (trazo discontinuo, forma
de Morse construida a partir de los valores medidos de $D_0$, $\omega_e$, $\omega_ex_e$, $r_e$). Cerca del
mínimo, RHF es bueno; a $R$ grande sube muy por encima de dos átomos de hidrógeno.}
\end{omfigure}

\begin{proposition}[Por qué el Hartree--Fock restringido falla en la disociación]\label{prop:b3:computational-chemistry:rhf-dissociation}
En la \omterm{def:b3:computational-chemistry:basis}{base mínima}, la función de onda RHF de \ce{H2} es, salvo
normalización, $\sigma_g(1)\sigma_g(2) \propto \chi_A(1)\chi_B(2) + \chi_B(1)\chi_A(2)
+ \chi_A(1)\chi_A(2) + \chi_B(1)\chi_B(2)$: a cualquier distancia, mitad covalente y mitad
iónica (\ce{H+ H-}). A $R$ grande, su energía tiende a la media de dos átomos
neutros y de un par de iones, no a la de dos átomos.
\end{proposition}

\begin{proof}
Al desarrollar $(\chi_A + \chi_B)(1)(\chi_A + \chi_B)(2)$, los dos términos cruzados ponen un
electrón en cada átomo y los dos términos cuadrados ambos en el mismo átomo, con
pesos iguales sea cual sea $R$. El determinante no puede cambiar esos pesos.
\end{proof}

\begin{definition}[Correlación electrónica, energía de correlación]\label{def:b3:computational-chemistry:correlation}
La \emph{correlación electrónica}\index{correlacion electronica@correlación electrónica} es la parte de la
evitación mutua de los electrones que no describe un único determinante (en el que los
electrones de espines opuestos se mueven independientemente, cada uno en el campo medio
del otro). La \emph{energía de correlación}\index{energia de correlacion@energía de correlación} es la energía
exacta menos la energía de Hartree--Fock en una base completa; siempre es
negativa.
\end{definition}

La \omterm{def:b3:computational-chemistry:correlation}{energía de correlación} es de unos \qty{-0.04}{\hartree} por par de electrones, un
uno por ciento de la energía total pero comparable a la energía de una reacción:
los métodos que la añaden (interacción de configuraciones, teoría de perturbaciones,
clústeres acoplados) cuestan mucho más que Hartree--Fock, cuyo coste ya crece
aproximadamente como la cuarta potencia del número de funciones de base.

\begin{theorem}[Teorema de Koopmans]\label{thm:b3:computational-chemistry:koopmans}
En la teoría de Hartree--Fock, la energía necesaria para arrancar un electrón del
orbital ocupado $a$, manteniendo congelados todos los demás orbitales, es $-\varepsilon_a$. Las
energías de ionización medidas son, por tanto, aproximadamente las opuestas de
las energías orbitales: es el \emph{teorema de Koopmans}\index{teorema de Koopmans}.
\end{theorem}

\begin{proof}
La energía de un determinante de capa cerrada es $\sum_a2h_{aa} + \sum_{ab}(2J_{ab} -
K_{ab})$, y $\varepsilon_a = h_{aa} + \sum_b(2J_{ab} - K_{ab})$. Se arranca un electrón
de $a$ sin cambiar los orbitales: los términos perdidos son $h_{aa}$, las
interacciones de ese electrón con todos los demás, $\sum_b(2J_{ab} - K_{ab}) -
J_{aa}$, más $J_{aa}$ con su antiguo compañero ---en total, exactamente
$\varepsilon_a$. Así, $E^+ - E = -\varepsilon_a$.
\end{proof}

\begin{example}[Ionización de \ce{H2}]\label{ex:b3:computational-chemistry:koopmans}
% ledger: ie4:H2
A \qty{1.4}{\bohr}, la energía orbital STO-3G de $\sigma_g$ es
\qty{-0.5782}{\hartree}: Koopmans predice \qty{15.73}{eV}, frente a los
\qty{15.43}{eV} medidos. Los orbitales congelados sobrestiman la energía del ion (que
se relajaría), y la correlación ausente actúa en sentido contrario; en las ionizaciones de
valencia los dos errores se compensan en gran parte, y por eso los espectros fotoelectrónicos
del volumen del segundo año pueden leerse con diagramas de orbitales.
\end{example}

\section{Teoría del funcional de la densidad, y lo que puede decir un cálculo}

La función de onda de $N$ electrones depende de $3N$ coordenadas; la
\omterm{def:b3:computational-chemistry:dft}{densidad electrónica} depende de tres.

\begin{definition}[Densidad electrónica, teoría del funcional de la densidad]\label{def:b3:computational-chemistry:dft}
La \emph{densidad electrónica}\index{densidad electronica@densidad electrónica} $\rho(\mathbf r)$ es el
número de electrones por unidad de volumen en $\mathbf r$, sumado sobre todos los electrones:
$\int\rho\,\dd\tau = N$. La \emph{teoría del funcional de la densidad}\index{teoria del funcional de la densidad@teoría del funcional de la densidad}
(DFT) calcula la energía del estado fundamental como un funcional de $\rho$. En la
práctica, $\rho = \sum_a|\phi_a|^2$ se construye a partir de \emph{orbitales de Kohn--Sham}\index{orbitales de Kohn--Sham},
los orbitales de unos electrones independientes ficticios
que tienen la misma densidad que los reales, y todos los efectos
de intercambio y de correlación se reúnen en un \emph{funcional de intercambio--correlación}\index{funcional de intercambio--correlacion@funcional de intercambio--correlación}
$E_{\mathrm{xc}}[\rho]$.
\end{definition}

\begin{theorem}[Hohenberg--Kohn]\label{thm:b3:computational-chemistry:hohenberg-kohn}
La densidad del estado fundamental de un sistema de electrones determina el potencial
externo (los núcleos) salvo una constante, y por tanto el hamiltoniano y todas las
propiedades del estado fundamental; y el funcional de la energía de la densidad es
mínimo en la densidad verdadera.
\end{theorem}
\admitted

La demostración corresponde a cursos más avanzados. Garantiza que existe un funcional
exacto, no cuál es: todo funcional práctico es una
aproximación, elegida dentro de una jerarquía ---densidad local, correcciones de
gradiente, híbridos que incorporan una parte de intercambio de Hartree--Fock--- y
calibrada sobre datos de referencia. Las ecuaciones de Kohn--Sham se resuelven con el
mismo bucle SCF que Hartree--Fock, con un coste similar, pero incluyen la correlación;
por eso la DFT se ha convertido en el caballo de batalla de la química.

\begin{method}[Leer un cálculo]\label{met:b3:computational-chemistry:read-calculation}
\begin{enumerate}
  \item Comprobar que el SCF ha convergido y que la \omterm{def:b3:computational-chemistry:pes}{optimización de geometría} ha terminado
        con fuerzas por debajo del umbral.
  \item Hacer un cálculo de frecuencias al mismo nivel: todas las frecuencias reales
        indican un mínimo; una frecuencia imaginaria, una estructura de transición, cuyo
        modo normal es el movimiento a través de la barrera.
  \item Comparar lo comparable: diferencias de energía entre estructuras
        calculadas a un mismo nivel, nunca energías absolutas de niveles distintos.
  \item Comparar con el experimento conociendo los errores típicos del nivel:
        las frecuencias armónicas suelen escalarse a la baja por 0.9 a 0.97;
        las longitudes de enlace son buenas a 1 pm aproximadamente con un funcional híbrido y una
        base polarizada; las energías de reacción, a unos pocos \unit{kJ/mol} en el mejor de los casos.
\end{enumerate}
\end{method}

\begin{method}[Elegir un nivel de teoría]\label{met:b3:computational-chemistry:choose-level}
Geometrías y frecuencias de moléculas orgánicas: un funcional híbrido con una
\omterm{def:b3:computational-chemistry:split-valence}{base de valencia desdoblada} polarizada. Interacciones débiles: añadir una corrección de
dispersión y \omterm{def:b3:computational-chemistry:split-valence}{funciones difusas}. Barreras de reacción y energética
precisa: métodos de función de onda correlacionados sobre geometrías DFT, hasta donde lo
permita el presupuesto. Sistemas grandes (proteínas, superficies): una pequeña región tratada
cuánticamente dentro de un campo de fuerzas clásico. Probar siempre el nivel
elegido en una molécula emparentada cuya respuesta se conozca.
\end{method}

\begin{inthelab}[Un experimento computacional]
Un cálculo se lleva a cabo como un experimento y se registra en el cuaderno: el
programa y su versión, el método, la base, los umbrales de convergencia,
la estructura de partida, y los ficheros de salida, conservados como datos brutos. Una secuencia
típica es: construir la molécula, optimizar su geometría a un nivel modesto,
calcular las frecuencias para confirmar un mínimo y después calcular la energía a un
nivel más alto sobre esa geometría.
\end{inthelab}

\begin{history}[Una molécula conocida primero por el cálculo]
% ledger: hist4:HeHp.2019
\ce{HeH+} se obtuvo en el laboratorio en 1925. Sus propiedades se calcularon
con detalle a partir de la década de 1930, y se predijo que sería abundante en
algunos gases astrofísicos; pero su primera raya de rotación está en el infrarrojo
lejano, absorbida por la atmósfera. En 2019, un observatorio aerotransportado
la detectó en la nebulosa planetaria NGC 7027, a la frecuencia que los cálculos
y los espectros de laboratorio habían predicho.
\end{history}

\section{Ejercicios}

\begin{exercise}[$\star$]\label{exo:b3:computational-chemistry:1}
Para el átomo de hidrógeno, se toma la función de prueba $\eu^{-\alpha r^2}$. En unidades
atómicas, $E(\alpha) = \frac32\alpha - 2\sqrt{2\alpha/\pi}$. Encuéntrense el mejor $\alpha$ y
la energía correspondiente; compárese con $-\frac12E_h$.
\end{exercise}

\begin{exercise}[$\star$]\label{exo:b3:computational-chemistry:2}
¿Cuántas funciones de base usa un cálculo sobre el benceno en la base
STO-3G? ¿Y en 6-31G(d), con seis funciones $d$ cartesianas por carbono y dos funciones $s$
por hidrógeno?
\end{exercise}

\begin{exercise}[$\star$]\label{exo:b3:computational-chemistry:3}
Explíquese por qué \ce{H2} y \ce{D2} tienen la misma longitud de enlace y la misma \omterm{def:b3:quantum-model-systems:oscillator}{constante
de fuerza} pero números de onda de vibración distintos. ¿Qué razón cabe esperar?
\end{exercise}

\begin{exercise}[$\star$]\label{exo:b3:computational-chemistry:4}
Un cálculo de frecuencias sobre una estructura de fórmula \ce{C2H5F} da 17 frecuencias
reales y una frecuencia imaginaria, $487\iu$~\unit{cm^{-1}}. ¿Qué tipo de
\omterm{def:b3:computational-chemistry:pes}{punto estacionario} es? ¿Cuántas frecuencias cabía esperar en total?
\end{exercise}

\begin{exercise}[$\star\star$]\label{exo:b3:computational-chemistry:5}
% ledger: ie:He, ie4:He2
Con $E(\zeta) = \zeta^2 - \frac{27}{8}\zeta$, calcúlese la mejor energía variacional
del helio en \unit{\hartree} y en eV, y su error respecto del valor exacto
$\qty{-2.9034}{\hartree}$. ¿Qué fracción de la energía total representa el error?
\end{exercise}

\begin{exercise}[$\star\star$]\label{exo:b3:computational-chemistry:6}
Dos funciones de base dan $H_{11} = \qty{-1.0}{\hartree}$, $H_{22} =
\qty{-0.5}{\hartree}$, $H_{12} = \qty{-0.2}{\hartree}$, $S_{12} = 0.3$. Resuélvase la
ecuación secular. ¿Queda la raíz más baja por debajo de $H_{11}$?
\end{exercise}

\begin{exercise}[$\star\star$]\label{exo:b3:computational-chemistry:7}
Un cálculo de Hartree--Fock tarda 10 minutos para el benceno en STO-3G. Suponiendo
un coste proporcional a la cuarta potencia del número de funciones de base,
estímese el tiempo en 6-31G(d).
\end{exercise}

\begin{exercise}[$\star\star$]\label{exo:b3:computational-chemistry:8}
% ledger: ie4:H2
La energía orbital STO-3G de \ce{H2} a \qty{1.4}{\bohr} es
\qty{-0.5782}{\hartree}. Dese la estimación de Koopmans de la energía de ionización en eV
y compárese con \qty{15.43}{eV}. Nómbrense dos razones de la diferencia y el sentido
de cada una.
\end{exercise}

\begin{exercise}[$\star\star$]\label{exo:b3:computational-chemistry:9}
En STO-3G, la energía de un átomo de hidrógeno es \qty{-0.4666}{\hartree} y la energía
RHF de \ce{H2} a \qty{10}{\bohr} es \qty{-0.5960}{\hartree}. ¿Cuánto queda la energía RHF
por encima de la de dos átomos a esa distancia, en eV? Explíquese con los términos
iónicos.
\end{exercise}

\begin{exercise}[$\star\star\star$]\label{exo:b3:computational-chemistry:10}
Para $\phi = c_1\chi_1 + c_2\chi_2$ con funciones reales, escríbase $E(c_1,c_2)$ y
dedúzcanse las dos ecuaciones seculares imponiendo $\partial E/\partial c_1 =
\partial E/\partial c_2 = 0$.
\end{exercise}

\begin{exercise}[$\star\star\star$]\label{exo:b3:computational-chemistry:11}
% ledger: aw:H, diat:H2.we
La energía RHF/STO-3G de \ce{H2} vale $-1.116871$, $-1.117501$ y
$\qty{-1.116714}{\hartree}$ a $R = 1.30$, 1.35 y \qty{1.40}{\bohr}. Estímese
la \omterm{def:b3:quantum-model-systems:oscillator}{constante de fuerza} por diferencias finitas, conviértase a \unit{N/m}
($\qty{1}{\hartree\per\bohr\squared} = \qty{1556.9}{N/m}$) y calcúlese el
número de onda armónico. Compárese con $\tilde\omega_e = \qty{4401}{cm^{-1}}$.
\end{exercise}

\begin{exercise}[$\star\star\star$]\label{exo:b3:computational-chemistry:12}
Debe calcularse la diferencia de energía entre dos confórmeros de un azúcar,
que se sabe que difieren en unos pocos \unit{kJ/mol}, uno de ellos con un enlace de hidrógeno
intramolecular. Elíjanse y justifíquense un método, un \omterm{def:b3:computational-chemistry:basis}{conjunto de base} y las comprobaciones
que se harían.
\end{exercise}

\section{Problema: HeH\texorpdfstring{$^+$}{+}, la primera molécula}

\begin{problem}[{Problema de fin de semana --- un cálculo de Hartree--Fock en base
mínima del ion hidruro de helio: la base, la estimación de núcleo, el
orbital convergido y la unión de un protón al helio}]
\label{pb:b3:computational-chemistry:1}
% ledger: bse:He.sto3g.a1, bse:He.sto3g.a2, bse:He.sto3g.a3, bse:H.sto3g.a1, pa:He, ie:He, ie4:He2
Se calcula \ce{HeH+} a $R = \qty{1.4632}{\bohr}$ en la base STO-3G (exponentes de He
6.3624, 1.1589, 0.31365; exponentes de H 3.4253, 0.62391, 0.16886; los
mismos tres coeficientes). La función 1 está sobre He, la función 2 sobre H. El programa
da (unidades atómicas)
\[
\mathbf S = \begin{pmatrix}1 & 0.5368\\ 0.5368 & 1\end{pmatrix},\quad
\mathbf h = \begin{pmatrix}-2.5983 & -1.4318\\ -1.4318 & -1.7318\end{pmatrix},\quad
\mathbf F_{\mathrm{final}} = \begin{pmatrix}-1.5902 & -1.0610\\ -1.0610 & -0.8340
\end{pmatrix},
\]
el orbital ocupado $\phi = 0.8766\chi_1 + 0.2025\chi_2$, las energías orbitales
$-1.6328$ y \qty{-0.1725}{\hartree}, y las energías SCF $-2.7978$, $-2.8404$,
$-2.8418$, $-2.8418$. Un átomo de helio en la misma base tiene
\qty{-2.8078}{\hartree}. $\qty{1}{\hartree} = \qty{27.211}{eV} =
\qty{2625.5}{kJ/mol}$.

\medskip\noindent\textbf{Parte I --- La base.}
\begin{enumerate}
  \item ¿Qué significa ``STO-3G''?
  \item ¿Por qué los exponentes del helio son mayores que los del hidrógeno? Compruébese
        que su razón es $(1.69/1.24)^2$.
  \item ¿Cuántas funciones de base, electrones y orbitales ocupados hay?
  \item Calcúlese la energía de repulsión nuclear.
  \item Interprétese $S_{12} = 0.5368$.
  \item ¿Por qué $h_{11}$ es más bajo que $h_{22}$?
\end{enumerate}

\medskip\noindent\textbf{Parte II --- La estimación de núcleo.}
\begin{enumerate}[resume]
  \item Escríbase la ecuación secular $\det(\mathbf h - \varepsilon\mathbf S) = 0$ como
        una ecuación de segundo grado en $\varepsilon$.
  \item Resuélvase.
  \item ¿Por qué no es más que un punto de partida?
  \item Dese la fórmula de la matriz de Fock en función de $\mathbf h$, de la
        matriz de densidad y de las integrales bielectrónicas.
  \item ¿Qué interacciones entre electrones contiene $\mathbf F$ que le faltan a $\mathbf h$?
  \item Explíquese por qué el procedimiento debe iterarse.
\end{enumerate}

\medskip\noindent\textbf{Parte III --- El orbital convergido.}
\begin{enumerate}[resume]
  \item Compruébese que $\phi$ está normalizado.
  \item Las poblaciones son los elementos diagonales de $P\mathbf S$, $2c_1^2 + 2c_1c_2S$
        sobre He y $2c_2^2 + 2c_1c_2S$ sobre H. Calcúlense.
  \item ¿Dónde está la carga positiva? ¿Se describe mejor la molécula
        como \ce{He + H+} o como \ce{He+ + H}?
  \item Úsese el \omterm{thm:b3:computational-chemistry:koopmans}{teorema de Koopmans} para estimar la energía necesaria para arrancar un
        electrón de \ce{HeH+}.
  \item Calcúlese la energía electrónica $E_{\mathrm{el}} = E - V_{\mathrm{nn}}$.
  \item ¿Cuántos ciclos fueron necesarios para converger a \qty{e-4}{\hartree}? ¿Por qué
        disminuyó la energía en cada ciclo?
  \item Con los exponentes del helio escalados a $\zeta = 2.0925$ en lugar de 1.69, el
        mismo programa da \qty{-2.8607}{\hartree}. ¿Qué base es mejor, y
        por qué puede afirmarse?
  \item ¿Qué representa la energía del orbital virtual, \qty{-0.1725}{\hartree}?
\end{enumerate}

\medskip\noindent\textbf{Parte IV --- La unión de un protón.}
\begin{enumerate}[resume]
  \item ¿Cuál es la energía de un protón desnudo? ¿Y la de \ce{He + H+} muy alejados, en esta
        base?
  \item Calcúlese la energía liberada por \ce{He + H+ -> HeH+} a este nivel, en
        \unit{kJ/mol}.
  \item La afinidad protónica medida del helio es \qty{177.8}{kJ/mol}.
        Coméntese.
  \item La energía exacta del helio es menos la suma de sus energías de
        ionización, 24.5874 y \qty{54.4178}{eV}. Calcúlese en \unit{\hartree}, así como
        el error del átomo de helio STO-3G.
  \item ¿Por qué un error de este tamaño puede dar aun así geometrías moleculares
        útiles?
  \item Enúnciese el resultado: la energía total RHF/STO-3G de \ce{HeH+} a
        \qty{1.4632}{\bohr}.
\end{enumerate}
\end{problem}
